\section{Introduction}
\label{sec:intro}

Measuring the cycle life of a lithium-ion cell means cycling it until its
capacity falls to 80\,\% of nominal, which takes months. Early prediction is
hard because capacity stays nearly flat until a late ``knee'', and because
nominally identical cells differ in life. Severson \textit{et
al.}~\cite{severson2019} predicted cycle life with about 9\,\% error from the
first 100 cycles using features of the discharge curve, chiefly the variance
of $\dq$, the change in the discharge capacity--voltage curve between cycles
10 and 100. This project reproduces that result and asks whether other
model families, features derived from degradation mechanisms or fewer cycles
do better, and how the models transfer to a batch tested differently:
\begin{enumerate}
  \item[\textbf{RQ1}] Using the public data, the paper's train/test split
  and its elastic-net models, do we obtain the test errors of its Table~1
  for feature sets F1--F3 (Table~\ref{tab:featsets})?
  \item[\textbf{RQ2}] With 41 training cells, do non-linear models, a
  Gaussian process (GP) or an additive spline model (GAM), beat the elastic
  net, (a)~for new cells from the training batches and (b)~for a later batch
  tested with a different procedure?
  \item[\textbf{RQ3}] Do features derived from degradation mechanisms
  (dQ/dV peaks, resistance, temperature, anode overhang, capacity budget),
  screened on training cells only, add information beyond $\dq$ and the
  paper's features?
  \item[\textbf{RQ4}] How does the error change when only the first
  $N\in\{20,\dots,150\}$ cycles are used, and what is the smallest useful
  $N$?
  \item[\textbf{RQ5}] Trained only on the three original batches, how
  accurate are the models on an unseen batch with different rests and charge
  termination, and which features cause their errors (SHAP)?
\end{enumerate}

Table~\ref{tab:featsets} defines the 12 features of the paper's sets
F1--F3. The $\Delta Q$ features correlate strongly with life both overall
and within each batch. The temperature integral correlates strongly overall
($-0.70$) but with the opposite sign within each batch, so it mainly
separates the batches; charge time helps only within b1, the one batch in
which it varies.

\begin{table*}[t]
\centering
\caption{The 12 features, all computed from cycles 2--100. $\Delta Q(V)=\dq$
on 1000 voltages from 3.5 to 2.0\,V of the 4\,C discharge; $Q_{d,n}$:
discharge capacity of cycle $n$; $\log=\log_{10}$. Sets: membership in
F1--F3. $r$: Pearson correlation with $\log_{10}$ life on the 41 training
cells (temperature integral: 40, one probe fault), overall and within
batches b1 and b2.}
\label{tab:featsets}
\small
\setlength{\tabcolsep}{3.5pt}
\begin{tabular}{@{}l>{\raggedright\arraybackslash}p{0.43\textwidth}lcccc@{}}
\toprule
Feature & Definition (how it is computed) & Unit & Sets & $r$ & $r_\text{b1}$ & $r_\text{b2}$ \\
\midrule
\multicolumn{7}{@{}l}{\textit{Discharge curve, cycle 100 vs.\ 10}~\cite{severson2019}} \\
$\log\mathrm{Var}\,\Delta Q$ & $\log|\mathrm{Var}_V\,\Delta Q(V)|$, spread of the curve change over voltage & -- & F1 F2 F3 & $-0.90$ & $-0.93$ & $-0.76$ \\
$\log|\min\Delta Q|$ & $\log|\min_V \Delta Q(V)|$, largest local change (on the plateau) & -- & F2 F3 & $-0.89$ & $-0.91$ & $-0.76$ \\
$\log|\mathrm{skew}\,\Delta Q|$ & $\log|$skewness of $\Delta Q(V)$ over $V|$ & -- & F2 & $+0.38$ & $+0.49$ & $-0.22$ \\
$\log|\mathrm{kurt}\,\Delta Q|$ & $\log|$kurtosis of $\Delta Q(V)$ over $V|$ & -- & F2 & $+0.20$ & $+0.63$ & $-0.45$ \\
\midrule
\multicolumn{7}{@{}l}{\textit{Capacity}~\cite{severson2019}} \\
$Q_{d,2}$ & discharge capacity of cycle 2 & Ah & F2 F3 & $+0.33$ & $-0.01$ & $+0.15$ \\
$\max Q_d - Q_{d,2}$ & highest capacity in cycles 2--100 minus $Q_{d,2}$ & Ah & F2 & $+0.15$ & $+0.54$ & $+0.16$ \\
Fade slope & slope $a$ of the least-squares line $Q_{d,n}\approx a\,n+b$, $n=2\dots100$ & Ah/cycle & F3 & $+0.40$ & $+0.46$ & $+0.59$ \\
Fade intercept & intercept $b$ of the same line & Ah & F3 & $+0.28$ & $+0.03$ & $+0.04$ \\
\midrule
\multicolumn{7}{@{}l}{\textit{Other measurements}~\cite{severson2019}} \\
Charge time & mean time from 0 to 80\,\% SOC, cycles 2--6 & min & F3 & $+0.77$ & $+0.77$ & $-0.16$ \\
Temp.\ integral & $\sum_{n=2}^{100}\int T\,dt$ over each cycle (gaps $>10$\,min left out) & $^\circ$C\,min & F3 & $-0.70$ & $+0.06$ & $+0.45$ \\
Min.\ resistance & lowest internal resistance in cycles 2--100 & $\Omega$ & F3 & $-0.39$ & $+0.27$ & $+0.12$ \\
Resistance change & $R_{100}-R_2$ & $\Omega$ & F3 & $+0.38$ & $-0.02$ & $-0.10$ \\
\bottomrule
\end{tabular}
\end{table*}
